\documentclass[11pt,a4paper]{article}
\usepackage[utf8]{inputenc}
\usepackage[margin=1in]{geometry}
\usepackage{booktabs}
\usepackage{amsmath}
\usepackage{amssymb}
\usepackage{hyperref}
\usepackage{graphicx}
\usepackage{microtype}
\usepackage{cite}

\title{\textbf{Cross-Lingual Scientific Evidence Retrieval:\\Measuring the Hindi--English Retrieval Gap on SciFact}}
\author{
  \textbf{Student Research Project} \\
  Department of Information Technology \\
  Indian Institute of Information Technology, Allahabad \\
  \texttt{Evaluating Multilingual Reliability of LLMs}
}
\date{October 2026}

\begin{document}

\maketitle

\begin{abstract}
Evaluating the multilingual reliability of large language models (LLMs) and information retrieval systems in scientific literature analysis requires testing cross-lingual evidence access. We investigate how retrieval effectiveness changes when scientific claims are formulated in Hindi versus English against an English-language scientific evidence corpus. Using the BEIR SciFact benchmark (5,183 scientific abstracts, 300 test claims), we evaluate three distinct retrieval paradigms: (1) BM25 lexical retrieval, (2) an English-oriented dense sentence encoder (\texttt{all-MiniLM-L6-v2}), and (3) a multilingual dense sentence encoder (\texttt{paraphrase-multilingual-MiniLM-L12-v2}). We measure Recall@1, Recall@5, Recall@10, and MRR@10, and analyze the cross-lingual retrieval gap ($\Delta = \text{Score}_{\text{English}} - \text{Score}_{\text{Hindi}}$). Our empirical results across 1,800 query-level evaluations demonstrate that BM25 experiences catastrophic degradation under translated Hindi claims (Recall@5 drops from 0.6969 in English to 0.0850 in Hindi; $\Delta = +0.6119, p < 10^{-10}$), surviving only on untranslated Latin acronyms and numbers. The English dense encoder exhibits near-total collapse on Hindi queries (Recall@5 = 0.0250; $\Delta = +0.7129$). In contrast, multilingual dense retrieval significantly outperforms both baselines on Hindi queries (Recall@5 = 0.1758; MRR@10 = 0.1333, $p < 10^{-6}$) and narrows the gap to $\Delta = +0.3542$. However, multilingual dense retrieval incurs a notable capacity tax on English queries (Recall@5 = 0.5300 vs. 0.7379 for English dense), and its absolute cross-lingual retrieval success remains severely bottlenecked. We provide detailed error categorization across 1,030 failure instances, identify biomedical translation artifacts, and document full reproduction artifacts.
\end{abstract}

\section{Introduction}
Large language models (LLMs) are increasingly integrated into scientific literature review, evidence synthesis, and claim verification pipelines \cite{wadden2020fact,thakur2021beir}. However, scientific literature analysis remains heavily English-centric. Researchers, students, and practitioners worldwide often formulate technical queries in non-English languages, such as Hindi, while the underlying body of peer-reviewed literature is almost exclusively published in English.

Before an LLM can reason over scientific claims, determine factual entailment, or generate grounded citations, an evidence retrieval component must locate relevant source abstracts. If retrieval fails, downstream verification inevitably fails or hallucinates. Evaluating cross-lingual scientific evidence retrieval is therefore a foundational prerequisite for assessing the multilingual reliability of LLM-based literature analysis systems.

This investigation addresses two core research questions:
\begin{itemize}
    \item \textbf{RQ1:} How does scientific evidence retrieval effectiveness change when English scientific claims are replaced with Hindi translations against an English scientific corpus?
    \item \textbf{RQ2:} Does multilingual dense retrieval reduce the Hindi--English retrieval gap compared with lexical retrieval and an English-oriented dense retriever?
\end{itemize}

\section{Related Work}
\subsection{Scientific Claim Verification and Retrieval}
The SciFact benchmark \cite{wadden2020fact} was introduced to evaluate automated scientific claim verification. SciFact claims are expert-written biomedical assertions paired with peer-reviewed abstracts and sentence-level evidence rationales. Thakur et al. \cite{thakur2021beir} incorporated SciFact into the BEIR benchmark for zero-shot information retrieval evaluation. Monolingual English benchmarks, however, do not evaluate cross-lingual access.

\subsection{Dense and Lexical Retrieval}
BM25 \cite{robertson2009bm25} represents the standard probabilistic lexical matching baseline, relying on exact term matching. In cross-lingual settings with different scripts (e.g., Devanagari vs. Latin), lexical overlap drops to zero unless technical terms or loanwords remain untranslated. Dense retrieval using Sentence-BERT \cite{reimers2019sentence} and multilingual knowledge distillation \cite{reimers2020multilingual} maps multilingual inputs into a shared semantic vector space, enabling cross-lingual nearest-neighbor search. MiniLM \cite{wang2020minilm} provides computationally efficient distilled architectures suitable for resource-constrained evaluation.

\section{Methodology}

\subsection{Dataset and Corpus}
We utilize the official BEIR release of SciFact \cite{thakur2021beir}:
\begin{itemize}
    \item \textbf{Corpus:} 5,183 scientific abstracts from PubMed/bioRxiv in English.
    \item \textbf{Queries:} 300 test scientific claims.
    \item \textbf{Relevance Judgments (qrels):} 339 query-document relevance annotations in the test split.
\end{itemize}

\subsection{Hindi Query Generation and Provenance}
To evaluate cross-lingual retrieval under controlled conditions without confounding corpus-side translation artifacts, we preserve the original English corpus and translate the 300 test claims into Hindi. Translations were generated via machine translation using MyMemory (89 queries) and local \texttt{Helsinki-NLP/opus-mt-en-hi} MarianMT (211 queries). Every row in \texttt{data/hindi\_queries.csv} maintains strict identifier alignment (\texttt{query\_id}), recording the source English text, Hindi text, translation method, and an explicit verification flag (\texttt{human\_checked=false}). We disclose that these translations constitute a translated-query stress test rather than naturally authored Hindi scientific prose.

\subsection{Retrieval Models}
We evaluate three candidate retrieval systems:
\begin{enumerate}
    \item \textbf{BM25 (Lexical):} Transparent BM25Okapi implementation with whitespace and punctuation tokenization.
    \item \textbf{EN-Dense (Monolingual Control):} \texttt{sentence-transformers/all-MiniLM-L6-v2} (384-dimensional embeddings, 6 layers, 22.7M parameters), trained primarily on English web text.
    \item \textbf{Multi-Dense (Cross-Lingual):} \texttt{sentence-transformers/paraphrase-multilingual-MiniLM-L12-v2} (384-dimensional embeddings, 12 layers, 118M parameters), aligned across 50+ languages including Hindi via multilingual distillation.
\end{enumerate}

All dense embeddings are $L_2$-normalized, and ranking is computed via inner product (equivalent to cosine similarity).

\subsection{Evaluation Metrics}
Following our frozen protocol (\texttt{configs/experiment.yaml}):
\begin{itemize}
    \item \textbf{Primary Metric:} Recall@5, measuring the proportion of gold relevant abstracts retrieved within top 5 candidates.
    \item \textbf{Secondary Metrics:} Recall@1, Recall@10, and Mean Reciprocal Rank (MRR@10).
    \item \textbf{Performance Gap:} $\Delta = \text{Metric}_{\text{English}} - \text{Metric}_{\text{Hindi}}$.
    \item \textbf{Statistical Testing:} Paired two-tailed $t$-tests across 300 query-level score differences.
\end{itemize}

\section{Experimental Results}
\label{sec:results}

Experiments were conducted on Windows 11 under Python 3.14.2 on CPU. Results are derived from 1,800 query-level runs and verified independently with 100\% metric agreement.

\begin{table}[ht]
\centering
\small
\caption{Aggregate Retrieval Metrics on BEIR SciFact (Test Split, $N = 300$ queries).}
\label{tab:metrics}
\begin{tabular}{llcccc}
\toprule
\textbf{Language} & \textbf{Method} & \textbf{Recall@1} & \textbf{Recall@5} & \textbf{Recall@10} & \textbf{MRR@10} \\
\midrule
English & BM25 & 0.4917 & 0.6969 & 0.7542 & 0.5945 \\
English & EN-Dense (\texttt{all-MiniLM-L6-v2}) & 0.4823 & \textbf{0.7379} & \textbf{0.7833} & \textbf{0.6047} \\
English & Multi-Dense (\texttt{paraphrase-multilingual}) & 0.3669 & 0.5300 & 0.6156 & 0.4550 \\
\midrule
Hindi & BM25 & 0.0617 & 0.0850 & 0.0950 & 0.0744 \\
Hindi & EN-Dense (\texttt{all-MiniLM-L6-v2}) & 0.0133 & 0.0250 & 0.0417 & 0.0206 \\
Hindi & Multi-Dense (\texttt{paraphrase-multilingual}) & \textbf{0.0892} & \textbf{0.1758} & \textbf{0.2208} & \textbf{0.1333} \\
\bottomrule
\end{tabular}
\end{table}

\begin{table}[ht]
\centering
\small
\caption{Hindi--English Performance Gap ($\Delta = \text{Score}_{\text{English}} - \text{Score}_{\text{Hindi}}$) with Paired Significance.}
\label{tab:gap}
\begin{tabular}{lccccc}
\toprule
\textbf{Method} & $\Delta\textbf{Recall@1}$ & $\Delta\textbf{Recall@5}$ & $\Delta\textbf{Recall@10}$ & $\Delta\textbf{MRR@10}$ & \textbf{Paired $t$-stat (R@5)} \\
\midrule
BM25 & +0.4300 & +0.6119 $\pm$ 0.0541 & +0.6592 & +0.5200 & $t = 22.18, p < 10^{-10}$ \\
EN-Dense & +0.4690 & +0.7129 $\pm$ 0.0507 & +0.7417 & +0.5842 & $t = 27.55, p < 10^{-10}$ \\
Multi-Dense & \textbf{+0.2778} & \textbf{+0.3542 $\pm$ 0.0563} & \textbf{+0.3947} & \textbf{+0.3217} & $t = 12.33, p < 10^{-10}$ \\
\bottomrule
\end{tabular}
\end{table}

As summarized in Table \ref{tab:metrics} and Table \ref{tab:gap}:
\begin{enumerate}
    \item \textbf{Impact of Translation (RQ1):} English scientific claims experience massive retrieval degradation when translated into Hindi. BM25 suffers an absolute Recall@5 loss of 61.19\% ($p < 10^{-10}$), and EN-Dense drops by 71.29\% ($p < 10^{-10}$).
    \item \textbf{Multilingual Mitigation (RQ2):} Multilingual dense retrieval significantly outperforms English dense on Hindi queries ($+0.1508$ Recall@5, $t = 6.93, p < 10^{-6}$) and cuts the cross-lingual gap by roughly half ($\Delta_{\text{Recall@5}} = 0.3542$ vs. $0.7129$).
    \item \textbf{Capacity Trade-Off:} On English claims, Multi-Dense lags EN-Dense by 20.79\% in Recall@5 (0.5300 vs. 0.7379), reflecting the capacity cost of multilingual distillation in compact models.
\end{enumerate}

\section{Error Analysis and Qualitative Inspection}
\label{sec:errors}
Analysis of the 1,030 queries with $\text{RR@10} = 0$ reveals five distinct error modes:
\begin{enumerate}
    \item \textbf{Cross-Script Lexical Disconnect:} BM25 on Hindi retrieved non-zero relevant documents solely when claims contained untranslated Latin acronyms (e.g., ``PrP'', ``UK'') or numbers (e.g., ``1/2000'').
    \item \textbf{Polysemy and Biomedical Terminology Mismatch:} General-domain translation tools misinterpret polysemous scientific terms. In Query 1 (\textit{``0-dimensional biomaterials show inductive properties''}), \textit{``inductive''} was rendered as \textit{``आगमनात्मक''} (formal logical induction) rather than biological osteoinduction, invalidating dense semantic matching.
    \item \textbf{Translation Malapropisms:} In Query 75, the scientific assertion \textit{``Active H. pylori urease has a polymeric structure that compromises two subunits...''} was translated literally as \textit{``समझौता करती है''} (makes a compromise/concession), disrupting semantic vector proximity.
    \item \textbf{Sub-space Misalignment:} Technical biomedical entities in Devanagari script mapped to semantic regions distant from their English counterpart abstracts.
    \item \textbf{Multilingual Model Under-capacity:} The 118M parameter multilingual model exhibited failure on 36.7\% of English queries (vs. 20.7\% for EN-Dense), showing that domain-specific expressiveness is diluted.
\end{enumerate}

\section{Threats to Validity and Limitations}
\begin{enumerate}
    \item \textbf{Machine Translation Artifacts:} Hindi queries were generated via automated machine translation without expert bilingual post-editing (\texttt{human\_checked=false}). The findings characterize retrieval under translated technical queries rather than naturally authored Hindi scientific discourse.
    \item \textbf{Domain and Language Scope:} Experiments are restricted to biomedical literature in SciFact and the Hindi--English language pair.
    \item \textbf{Retrieval vs. Verification:} Recall@k measures evidence retrieval, not claim entailment or LLM factual reasoning.
\end{enumerate}

\section{Conclusion}
This study quantifies the cross-lingual scientific evidence retrieval gap between Hindi claims and an English scientific corpus on SciFact. While multilingual dense retrieval significantly outperforms lexical and English-oriented baselines on Hindi queries and halves the performance gap, substantial performance penalties and multilingual capacity constraints remain. Developing specialized cross-lingual scientific retrieval adapters is essential for trustworthy multilingual LLM-assisted literature analysis.

\bibliographystyle{plain}
\bibliography{references}

\end{document}
